\heading{实验物理中的统计方法-24春-期中}
\noteinfo{内容精修版；已统一题干表述、边界条件与解答细节。}

\begin{question}
$X\sim N(0,2^2)$，$Y\sim N(0,2^2)$，且 $V[X-Y]=0$，求 $X$ 和 $Y$ 的协方差矩阵。
\begin{solution}
由
\[
0=\V(X-Y)=\V(X)+\V(Y)-2\cov(X,Y)
=4+4-2\cov(X,Y),
\]
得 $\cov(X,Y)=4$。因此
\ansmath{\Sigma=
\begin{pmatrix}
4&4\\ 4&4
\end{pmatrix}}
\end{solution}
\end{question}

\begin{question}
有 $n\ge2$ 个编号为 1 到 $n$ 的小球，$n$ 个编号为 1 到 $n$ 的盒子，将小球放入盒子之中，每个盒子放且仅放一个小球。小球编号与盒子编号相等称之为配对。求配对数 $X$ 的方差。
\begin{solution}
记 $I_i=\mathbf 1\{\pi(i)=i\}$，则 $X=\sum_{i=1}^n I_i$。有
\[
E[I_i]=\frac1n,
\qquad E[I_iI_j]=\frac{(n-2)!}{n!}=\frac{1}{n(n-1)}\ (i\ne j).
\]
故
\[
E[X]=n\cdot\frac1n=1,
\]
\[
E[X^2]=\sum_iE[I_i]+2\sum_{i<j}E[I_iI_j]
=1+2\binom n2\frac{1}{n(n-1)}=2.
\]
于是
\[
\V(X)=E[X^2]-E[X]^2=2-1=1.
\]
\end{solution}
\end{question}

\begin{question}
一个仪器在任意 $t$ 时间内发生故障的次数 $N_t$ 服从参数为 $\lambda t$ 的泊松分布，求第一次发生故障的时间 $T$ 的概率密度函数。
\begin{solution}
\[
P(T>t)=P(N_t=0)=e^{-\lambda t},\qquad t\ge 0.
\]
故
\[
F_T(t)=1-e^{-\lambda t},\qquad
f_T(t)=\lambda e^{-\lambda t}\ (t\ge 0).
\]
\end{solution}
\end{question}

\begin{question}
平面上有方格子，边长为 $a>1$。将直径为 1 的圆形硬币扔在平面上，求硬币不与方格线相交的概率。
\begin{solution}
硬币半径为 $1/2$。在每个边长为 $a$ 的小方格中，圆心必须离四条边都至少 $1/2$，可落区域面积为 $(a-1)^2$。故概率
\ansmath{P=\frac{(a-1)^2}{a^2}}
\end{solution}
\end{question}

\begin{question}
一种原子的半衰期为 10 年。现有 5 个原子，经过了 20 年，恰好衰变 2 个的概率是多少？
\begin{solution}
单个原子在 $20$ 年内衰变概率
\[
p=1-\left(\frac12\right)^2=\frac34.
\]
故所求为二项分布概率
\ansmath{P=\binom52\left(\frac34\right)^2\left(\frac14\right)^3=\frac{45}
{512}}
\end{solution}
\end{question}

\begin{question}
一个显示屏随机显示从 1 到 $N$ 的数字，显示每个数字的概率相等。一个人观测此显示屏，看到显示屏前后共显示了 $n$ 个，此人记录下观测到的最大数字 $k$。求 $k$ 的概率分布。
\begin{solution}
\[
P(K\le k)=\left(\frac{k}{N}\right)^n,\qquad k=1,2,\dots,N.
\]
故
\ansmath{P(K=k)=P(K\le k)-P(K\le k-1)=\frac{k^n-(k-1)^n}
{N^n}}
\end{solution}
\end{question}

\begin{question}
元件发生故障的时刻 $\tau$ 的累积分布函数为 $F(\tau)$。现有两个相互独立的元件，只有它们都正常工作时仪器才能正常工作，求仪器正常工作的时间 $T$ 的累积分布函数。
\begin{solution}
\[
T=\min(T_1,T_2),\qquad
P(T>t)=[1-F(t)]^2,
\]
故
\ansmath{F_T(t)=1-[1-F(t)]^2}

\end{solution}
\end{question}

\begin{question}
$X$ 和 $Y$ 相互独立且都服从标准正态分布，$Z=X-Y$，求 $E[Z]$，$E[|Z|]$，$V[|Z|]$。
\begin{solution}
有 $Z\sim N(0,2)$，所以
\[
E[Z]=0,
\qquad E[|Z|]=\sqrt2\sqrt{\frac{2}{\pi}}=\frac{2}{\sqrt\pi}.
\]
又
\[
E[|Z|^2]=E[Z^2]=2,
\]
故
\[
\V(|Z|)=2-\left(\frac{2}{\sqrt\pi}\right)^2=2-\frac{4}{\pi}.
\]
\end{solution}
\end{question}

\begin{question}[10 分]
雷达显示屏是一个半径为 $R$ 的圆形区域，圆内随机出现一个光点且其位置服从均匀分布。求光点与圆心距离的期望和方差。
\begin{solution}
距离 $r$ 的密度为
\[
f(r)=\frac{2r}{R^2},\qquad 0<r<R.
\]
因此
\[
E[r]=\int_0^R r\frac{2r}{R^2}dr=\frac{2R}{3},
\qquad
E[r^2]=\int_0^R r^2\frac{2r}{R^2}dr=\frac{R^2}{2}.
\]
故
\[
\V(r)=\frac{R^2}{2}-\left(\frac{2R}{3}\right)^2=\frac{R^2}{18}.
\]
\end{solution}
\end{question}

\begin{question}[10 分]
一个矩形，边长分别为 $a$ 和 $b$，且 $a>b$。在矩形的任意边上任取两个点，求这两个点的距离的平方的期望值。
\begin{solution}
设边界上一点坐标为 $(X,Y)$，则两点距离平方
\[
D^2=(X_1-X_2)^2+(Y_1-Y_2)^2,
\]
故
\[
E[D^2]=2\V(X)+2\V(Y).
\]
对边界均匀分布，可算得
\[
E[X]=\frac a2,
\qquad
E[X^2]=\frac{2\int_0^a x^2dx+ba^2}{2(a+b)},
\]
\[
E[Y]=\frac b2,
\qquad
E[Y^2]=\frac{2\int_0^b y^2dy+ab^2}{2(a+b)}.
\]
化简后
\ansmath{E[D^2]=\frac{a^2+2ab+b^2}{6}=\frac{(a+b)^2}{6}}
\end{solution}
\end{question}

\begin{question}[10 分]
粒子束流包含 $10^{-4}$ 的电子，其余为光子。粒子通过某双层探测器，可能在 2 层都给出信号，也可能只有一层给出信号或者没有任何信号。电子 $(e)$ 和光子 $(\gamma)$ 在穿过该双层探测器给出 0、1 或 2 个信号的概率如下
\[
P(0\mid e)=0.001,\qquad P(1\mid e)=0.01,\qquad P(2\mid e)=0.989,
\]
\[
P(0\mid \gamma)=0.99899,\qquad P(1\mid \gamma)=0.001,\qquad P(2\mid \gamma)=10^{-5}.
\]
\begin{enumerate}
\item 如果只有一层给出了信号，该粒子为光子的概率是多少？
\item 如果两层都给出了信号，该粒子为电子的概率是多少？
\end{enumerate}
\begin{solution}
由 Bayes 公式：
\begin{enumerate}
\item
\ansmath{P(\gamma|1)=\frac{P(1|\gamma)P(\gamma)}{P(1|\gamma)P(\gamma)+P(1|e)P(e)}
\approx 0.9990009}
\item
\ansmath{P(e|2)=\frac{P(2|e)P(e)}{P(2|e)P(e)+P(2|\gamma)P(\gamma)}
\approx 0.908181}
\end{enumerate}
\end{solution}
\end{question}

\begin{question}[10 分]
某一个衰变事件（具体是什么我不记得了）的发生是极为稀有的事情。本底数为 3，现在观察到 6 个疑似信号，求显著性水平。
\begin{solution}
\ansmath{p=P(N\ge 6\mid \lambda=3)
=1-\sum_{k=0}^{5}e^{-3}\frac{3^k}{k!}
\approx 0.08392}
\end{solution}
\end{question}

\begin{question}[10 分]
制造了一批工件，每个工件的横截面积（单位为平方毫米）服从 $N(\mu,0.02^2)$。现在抽取 15 个样本，测得样本方差为 $S^2=0.025^2$。已知
\[
\frac{(n-1)S^2}{\sigma^2}\sim \chi^2(n-1),
\]
问方差有没有明显差异？参考数据：卡方分布的两个上 $\alpha$ 分位数为 $\chi^2_{0.025}(14)=26.119$，$\chi^2_{0.975}(14)=5.629$。
\begin{solution}
设原假设 $H_0:\sigma^2=0.02^2$。统计量
\[
\chi^2=\frac{(n-1)S^2}{\sigma_0^2}=14\times\frac{0.025^2}{0.02^2}=21.875.
\]
给定双侧临界区间
\[
\chi^2_{0.975}(14)=5.629,\qquad \chi^2_{0.025}(14)=26.119.
\]
由于
\ansmath{5.629<21.875<26.119,}
故不拒绝 $H_0$，即方差与 $0.02^2$ \textbf{无显著差异}。

\end{solution}
\end{question}

\begin{question}[7 分]
蒲丰投针实验计算 $\pi$ 的公式为 $\pi\approx\dfrac{2lN}{an}$，其中 $N$ 为总投掷次数，$n$ 为针与线相交次数。某次实验中，$l=2.5\mathrm{cm}$，$a=3\mathrm{cm}$，$N=3408$，$n=1808$，于是算出 $\pi=3.1415929$。忽略 $l$ 和 $a$ 的误差，计算 $\pi$ 的不确定度。
\begin{solution}
把 $n$ 视为二项分布 $\mathrm{Bin}(N,p)$，其中 $p\approx n/N$。于是
\[
\sigma_n\approx \sqrt{Np(1-p)}=\sqrt{n\left(1-\frac{n}{N}\right)}\approx 29.13.
\]
由误差传播
\[
\sigma_{\pi}\approx \left|\frac{\partial \hat\pi}{\partial n}\right|\sigma_n
=\frac{2lN}{a n^2}\sigma_n=\hat\pi\frac{\sigma_n}{n}.
\]
代入
\ansmath{\hat\pi\approx 3.1415929,
\qquad \sigma_{\pi}\approx 0.0506}
故可写为 $\pi\approx 3.14\pm 0.05$。
\end{solution}
\end{question}

\begin{question}[6 分]
仪器在任意时间内观察到噪声的概率相等。在 $[0,T]$ 的时间窗口内观察到两个噪声，如果两个噪声间距小于 $t$，就被误鉴别为信号。在时间窗口内已经观察到两个噪声的前提下，求误鉴别为信号的概率。
\begin{solution}
条件于“恰有两个噪声”后，两到达时刻等价于两个独立均匀点 $U,V\sim U(0,T)$。所求为
\[
P(|U-V|<t)=1-\frac{2\cdot \frac12(T-t)^2}{T^2}
=1-\frac{(T-t)^2}{T^2}
=\frac{2t}{T}-\left(\frac{t}{T}\right)^2,
\]
其中 $0\le t\le T$；若 $t\ge T$，该概率为 $1$。
\ansmath{P(|U-V|<t)=1-\frac{2\cdot \frac12(T-t)^2}{T^2}
=1-\frac{(T-t)^2}{T^2}
=\frac{2t}{T}-\left(\frac{t}{T}\right)^2,}
\end{solution}
\end{question}

\begin{question}
已知 $X_i\sim U(0,1)$。定义随机变量
\[
N_a=\min\left\{k:\prod_{i=1}^{k}X_i<a\right\},\qquad 0<a<1.
\]
求此随机变量的概率分布。
\begin{solution}
令 $Y_i=-\ln X_i$，则 $Y_i\sim \mathrm{Exp}(1)$，且
\[
\prod_{i=1}^k X_i<a
\iff \sum_{i=1}^k Y_i>-\ln a.
\]
记 $\lambda=-\ln a$。则 $N_a=k$ 当且仅当
\[
S_{k-1}\le \lambda<S_k,
\]
其中 $S_k=\sum_{i=1}^k Y_i$ 是单位速率 Poisson 过程的到达时刻。故
\[
N_a-1\sim \mathrm{Poisson}(\lambda),
\]
即
\[
P(N_a=k)=e^{-\lambda}\frac{\lambda^{k-1}}{(k-1)!}
=a\frac{(-\ln a)^{k-1}}{(k-1)!},
\qquad k=1,2,\dots.
\]
\ansmath{P(N_a=k)=e^{-\lambda}\frac{\lambda^{k-1}}{(k-1)!}
=a\frac{(-\ln a)^{k-1}}{(k-1)!},
\qquad k=1,2,\dots}
\end{solution}
\end{question}


